\section*{Chapter \ref{ch:b3:climate-models} --- Climate Models}

\begin{solution}{exo:b3:climate-models:1}
$A = 14.0 - 1.5 = \num{12.5}$ and $B = 0.14 - 0.05 = \num{0.09}$, in
\unit{kcal.cm^{-2}} per month. A month is \qty{2.63e6}{s}, so
\qty{1}{kcal.cm^{-2}} per month is $\num{4.184e7}/\num{2.63e6} =
\qty{15.91}{W/m^2}$: $A = \qty{199}{W/m^2}$, $B =
\qty{1.43}{W.m^{-2}.\celsius^{-1}}$. At \qty{15}{\celsius}, $199 + 21.5 =
\qty{220}{W/m^2}$.
\end{solution}

\begin{solution}{exo:b3:climate-models:2}
$(238.0 - 211.9)/1.816 = \qty{14.4}{\celsius}$; with 0.31, $(234.6 -
211.9)/1.816 = \qty{12.5}{\celsius}$. Each 0.01 of albedo costs
\qty{3.4}{W/m^2} and $3.4/1.816 = \qty{1.9}{\celsius}$.
\end{solution}

\begin{solution}{exo:b3:climate-models:3}
$C = 1025\times3992\times60 = \qty{2.46e8}{J.m^{-2}.K^{-1}}$; divided by
the \qty{3.156e7}{s} of a year, \qty{7.78}{W.yr.m^{-2}.K^{-1}}. Then $\tau
= 7.78/1.816 = \qty{4.3}{years}$, and $7.78/1.33 = \qty{5.8}{years}$.
\end{solution}

\begin{solution}{exo:b3:climate-models:4}
ECS $= 3.93/1.33 = \qty{2.95}{\celsius}$; TCR $\approx 3.93/1.95 =
\qty{2.02}{\celsius}$; ratio \num{0.68}.
\end{solution}

\begin{solution}{exo:b3:climate-models:5}
$Q = \qty{335}{W/m^2}$. Snowball: $(335\times0.38 - 211.9)/1.816 =
\qty{-46.6}{\celsius}$; warm: $(335\times0.70 - 211.9)/1.816 =
\qty{12.4}{\celsius}$, still above \qty{10}{\celsius}; ramp: $335(0.54 +
0.016\,T) = 211.9 + 1.816\,T$ gives $T = \qty{8.7}{\celsius}$. $N' =
\qty{-1.816}{W.m^{-2}.\celsius^{-1}}$ at the first two, stable, and $-1.816
+ 335\times0.016 = \num{+3.54}$ at the ramp, unstable.
\end{solution}

\begin{solution}{exo:b3:climate-models:6}
The warm branch ends where its temperature reaches \qty{0}{\celsius}:
$0.70\,Q = A$, $Q = \qty{302.7}{W/m^2}$, a solar constant of
\qty{1211}{W/m^2}, \qty{11.0}{\percent} below today's. The snowball melts
where $0.38\,Q = A - 20B$, $Q = \qty{462.1}{W/m^2}$, $S_0 =
\qty{1848}{W/m^2}$. The threshold is $B(T^* - T_F)/(Q_0(1 - \alpha_F))$,
the gap between today's temperature $T^*$ and the temperature at which the
ice starts to spread, converted to sunlight: $1.816\times14.37/238 =
\qty{11.0}{\percent}$ here, $1.816\times4.37/238 = \qty{3.3}{\percent}$ in
the text.
\end{solution}

\begin{solution}{exo:b3:climate-models:7}
(a)~\qty{-45.5}{\celsius}. (b)~$340\times0.38 + F = 211.9 - 18.2$: $F =
\qty{64.5}{W/m^2}$, 16 times the forcing of a doubling. (c)~While the
forcing lasts the planet jumps to the warm branch, $(238.0 + 64.5 -
211.9)/1.816 = \qty{50}{\celsius}$; once it is removed the planet cools
along the warm branch to \qty{14.4}{\celsius}, which exists at today's Sun:
it does not return to the snowball.
\end{solution}

\begin{solution}{exo:b3:climate-models:8}
$\Sigma = 1.80/7.0 + 0.70/90 = \qty{0.2649}{yr^{-1}}$, $\Pi = 1.10\times0.70/
(7.0\times90) = \qty{1.222e-3}{yr^{-2}}$. The roots give $\tau_f =
\qty{3.84}{years}$ and $\tau_s = \qty{213}{years}$; the approximations,
$7.0/1.80 = \qty{3.89}{years}$ and $90\times1.80/(1.10\times0.70) =
\qty{210}{years}$. $a_f = 0.597$ exactly, $1.10/1.80 = 0.611$
approximately. ECS $= 3.93/1.10 = \qty{3.57}{\celsius}$; TCR $\approx
3.93/1.80 = \qty{2.18}{\celsius}$.
\end{solution}

\begin{solution}{exo:b3:climate-models:9}
Denominator $B + \beta = \num{5.555}$. Equator: $(340\times1.221\times0.70 -
211.9 + 3.74\times12.81)/5.555 = \qty{22.8}{\celsius}$; pole:
$(340\times0.506\times0.38 - 211.9 + 47.9)/5.555 = \qty{-17.7}{\celsius}$.
The transport removes $3.74\times(22.8 - 12.8) = \qty{37}{W/m^2}$ at the
equator and brings $3.74\times30.6 = \qty{114}{W/m^2}$ to the pole.
\end{solution}

\begin{solution}{exo:b3:climate-models:10}
(a)~$(1 - x^2)P_2' = 3x - 3x^3$, whose derivative is $3 - 9x^2 = -6P_2$.
(b)~The transport term is $-\mathcal D T_2(-6P_2) = 6\mathcal DT_2P_2$.
The balance $Q(1 + s_2P_2)(1 - \alpha) - A - B(T_0 + T_2P_2) - 6\mathcal
DT_2P_2 = 0$ holds for all $x$ if its constant part and its $P_2$ part
vanish, which gives $T_0$ and $T_2$. (c)~$T(0) - T(1) = T_2(-\tfrac12 - 1)
= -1.5\,T_2$. Without transport $T_2 = 340\times(-0.476)\times0.70/1.816 =
\qty{-62.4}{\celsius}$: a difference of \qty{94}{\celsius}. (d)~$T_2 =
-55.3/1.5 = \qty{-36.9}{\celsius}$, so $B + 6\mathcal D = 113.3/36.9 =
3.07$ and $\mathcal D = \qty{0.21}{W.m^{-2}.\celsius^{-1}}$.
\end{solution}

\begin{solution}{exo:b3:climate-models:11}
(a)~$r = \tfrac12\Sigma\big(1 \pm \sqrt{1 - 4\Pi/\Sigma^2}\big)$; with
$\sqrt{1 - \epsilon} \approx 1 - \epsilon/2$, the roots are $\Sigma -
\Pi/\Sigma \approx \Sigma$ and $\Pi/\Sigma$. (b)~With $T_d = 0$, $C\,\mathrm
dT/\mathrm dt = kt - (\gamma - \lambda)T$, whose solution from rest is $T =
k(t - \tau_f)/(\gamma - \lambda) + k\tau_f\mathrm e^{-t/\tau_f}/(\gamma -
\lambda)$, $\tau_f = C/(\gamma - \lambda)$. For $t \gg \tau_f$, $T \approx
kt/(\gamma - \lambda)$, and at the doubling, $kt = F_{2\times}$: TCR
$\approx F_{2\times}/(\gamma - \lambda)$. (c)~$\Sigma = 1.95/8.1 +
0.62/110 = \qty{0.2464}{yr^{-1}}$, $\Pi = 1.33\times0.62/(8.1\times110) =
\qty{9.25e-4}{yr^{-2}}$; $\tau_f = \qty{4.12}{years}$ and $\tau_s =
\qty{262}{years}$ (approximately $1/\Sigma = \qty{4.06}{years}$ and
$\Sigma/\Pi = \qty{266}{years}$); $4\Pi/\Sigma^2 = 0.061$.
\end{solution}

\begin{solution}{exo:b3:climate-models:12}
(a)~Mean \qty{3.50}{\celsius}, standard deviation \qty{0.81}{\celsius}.
(b)~ECS $= 1.88 + 1.86\,X$; correlation \num{0.992}; residual standard
deviation \qty{0.11}{\celsius}. (c)~$1.88 + 1.86\times0.7 =
\qty{3.18}{\celsius}$, with a spread $\sqrt{(1.86\times0.1)^2 + 0.11^2} =
\qty{0.22}{\celsius}$, about a quarter of the spread of the ensemble.
(d)~No physical mechanism linking $X$ to the feedbacks; an $X$ found by
trying many observables, so that the correlation may be chance; a bias
shared by all eight models; models of the same family counted as
independent.
\end{solution}

\begin{solution}{pb:b3:climate-models:1}
\textbf{1.} $Q_0 = 1360/4 = \qty{340}{W/m^2}$: the planet intercepts the
sunlight on a disc of area $\pi R^2$ and spreads it over a sphere of area
$4\pi R^2$.

\textbf{2.} $340\times0.70 = \qty{238.0}{W/m^2}$; $211.9 + 1.816\times15 =
\qty{239.1}{W/m^2}$.

\textbf{3.} $(238.0 - 211.9)/1.816 = \qty{14.4}{\celsius}$.

\textbf{4.} $(129.2 - 211.9)/1.816 = \qty{-45.5}{\celsius}$, below
\qty{-10}{\celsius}.

\textbf{5.} $340(0.54 + 0.016\,T) = 211.9 + 1.816\,T$: $T =
(183.6 - 211.9)/(1.816 - 5.44) = \qty{7.8}{\celsius}$.

\textbf{6.} $(238.0/\num{5.670e-8})^{1/4} = \qty{254.5}{K}$,
\qty{-18.6}{\celsius}: \qty{33}{\celsius} below the warm equilibrium, the
greenhouse effect of the planet built into $A$ and $B$.

\textbf{7.} $N' = -B = \qty{-1.816}{W.m^{-2}.\celsius^{-1}}$ at the
snowball and the warm equilibrium, both stable; $-1.816 + 5.44 =
\num{+3.62}$ at the middle one, unstable.

\textbf{8.} $\tau = 8.1/3.62 = \qty{2.2}{years}$.

\textbf{9.} $8.1/1.816 = \qty{4.5}{years}$.

\textbf{10.} $3.93/1.816 = \qty{2.2}{\celsius}$, below the assessed best
estimate of \qty{3}{\celsius} and below its \emph{likely} range, 2.5 to
\qty{4}{\celsius}: the $B$ of the model, fitted across latitudes, is not
the feedback parameter of the climate as a whole.

\textbf{11.} $\mathrm e^{-t/\tau} = 0.05$ at $t = 3\tau =
\qty{13}{years}$. The real climate also warms the deep ocean, whose time
scale is centuries.

\textbf{12.} At $T_F = \qty{10}{\celsius}$; $Q = (211.9 + 18.2)/0.70 =
\qty{328.7}{W/m^2}$.

\textbf{13.} $S_0 = \qty{1315}{W/m^2}$, \qty{3.3}{\percent} below today's.

\textbf{14.} To the snowball, $(328.7\times0.38 - 211.9)/1.816 =
\qty{-47.9}{\celsius}$.

\textbf{15.} $Q = (211.9 - 18.2)/0.38 = \qty{509.8}{W/m^2}$, $S_0 =
\qty{2039}{W/m^2}$, \qty{50}{\percent} above today's.

\textbf{16.} Down the warm branch to \qty{10}{\celsius} at
\qty{1315}{W/m^2}, a fall to \qty{-48}{\celsius}, then back up the
snowball branch to \qty{-45.5}{\celsius} at \qty{1360}{W/m^2}: the planet
stays frozen. Only beyond \qty{2039}{W/m^2} would it jump back to the warm
branch.

\textbf{17.} $\bar\alpha = 0.30\times0.9738 + 0.62\times0.0262 =
\num{0.3084}$; $\bar T = (340\times0.6916 - 211.9)/1.816 =
\qty{12.80}{\celsius}$.

\textbf{18.} $T_c = (340\times0.5643\times0.50 - 211.9 + 3.739\times12.80)/
5.555 = \qty{-12.26}{\celsius}$, \qty{3.5}{\celsius} warmer than the
reanalysis at \ang{72}.

\textbf{19.} $5.555\times(-12.26) + 211.9\times3.059 = \num{580.1}$, with
$\beta/B = \num{2.059}$.

\textbf{20.} The denominators are 1.7053, 1.7473, 1.7565, 1.7488 and
1.7362, and $Q/Q_0 = \num{1.0005}$, \num{0.9764}, \num{0.9713},
\num{0.9756} and \num{0.9827}.

\textbf{21.} Near \ang{47.5}, at \num{0.971}. The ice lines poleward of the
minimum are stable, $\mathrm dQ/\mathrm dx_s > 0$; those equatorward are
unstable.

\textbf{22.} $Q/Q_0 = 1$ also between \ang{30} and \ang{40}, close to
\ang{30} (\ang{30.2} by interpolation): an unstable ice line. At \ang{0},
$Q/Q_0 = 580.1/(0.6106 + 0.7824)/340 = \num{1.225}$: the snowball holds
for any Sun up to \qty{22}{\percent} brighter than today's, so it is
possible today.

\textbf{23.} \qty{2.9}{\percent}, against \qty{3.3}{\percent} for the
zero-dimensional model and \qty{1.6}{\percent} for Budyko; the ice stands
near \ang{48} at the threshold, against his \ang{50}. The latitudinal model
glaciates sooner than the global one, because the sunlight is weakest where
the ice forms.

\textbf{24.} The model planet glaciates for a fall of
\textbf{\qty[text-series-to-math=true]{2.9}{\percent}} of the solar
constant (\qty{2.87}{\percent} by the calculation).
\end{solution}
